\answer{ex:20:1}{$t_{\text{б}}=\tau_r\ln\frac{1-L}{1-H}=16.8$~ч, $t_{\text{с}}=\tau_d\ln\frac HL=5.8$~ч, период $22.5$~ч; к $24$~ч генератор подтягивает суточное качание порогов --- фазовая автоподстройка~\eqref{eq:pll}.}
\answer{ex:20:2}{(а)~$L=0.111$; при $L=0.092$ сон длился бы $9.6$~ч. (б)~$S_0$ больше в $1.44$ раза, сон длиннее в $1.10$ раза: скорость разряда пропорциональна $S$. (в)~На $22\%$ ($1/0.82=1.22$).}
\answer{ex:20:3}{$H=\frac{p-1}{p-1/q}=0.623$, $L=H/q=0.093$, где $p=e^{16/\tau_r}$, $q=e^{8/\tau_d}$. После побудки через $4$~ч фаза сдвигается навсегда (засыпание на $7.2$~ч раньше); с качающимися порогами она возвращается за двое--трое суток.}
\answer{ex:20:4}{(а)~$r_\pm=\frac12\bigl(1\pm\sqrt{1-4k/w}\bigr)=0.947,\ 0.053$; $a_{\text{в}}=3.51$, $a_{\text{н}}=0.49$. (б)~$T_{\text{раб}}\approx0.33$~с, $T_{\text{тиш}}\approx0.59$~с, $T\approx0.93$~с, доля работы $0.36$. (в)~Модель $1.11$~с из-за конечного $\tau/\tau_a$; при росте $\tau_a$ отношение $T/\tau_a$ стремится к $3.2$. Быстрее суточного в $7.8\cdot10^4$ раз.}
\answer{ex:20:5}{(а)~$a_{\text{в}}/r_+<b<a_{\text{н}}/r_-$: $3.71<b<9.20$. (б)~$r\approx1$, $a\approx1$, $F'\approx0$ --- устойчиво. (в)~\nt{АЦЕТ} уменьшает $b$ ниже $3.71$, и пересечение уходит на верхнюю ветвь.}
\answer{ex:20:6}{(а)~$16.79$ и $5.76$~ч. (б)~При $D=24$, $32$, $40$, $48$, $64$ сон длится $4.9$, $6.8$, $8.8$, $5.5$, $8.9$~ч и совпадает с $\tau_d\ln(S_0/L)$ до $0.01$~ч; длительность задаёт суточная фаза засыпания, а не долг. (в)~Колебания при $3.8\le b\le9.0$ (нет при $3.7$ и $9.2$); период $1.61$~с при $b=3.8$, $1.11$~с при $b=5$, минимум $0.98$~с около $b=8$.}
\answer{ex:20:7}{(а)~Ошибка на A после B в среднем $0.51$ (от $0.19$ до $1.08$), ошибка нулевого ответа --- $1$. (б)~Обе ошибки меньше $10^{-4}$. (в)~Изменения лежат в линейной оболочке образов B; общее решение есть, пока $P_A+P_B\le n$.}
\answer{ex:20:8}{Открытая задача. Равномерное масштабирование сохраняет отношения весов, но сохраняет ответы порогового нейрона лишь при масштабировании порога в то же число раз; перемешанные повторы меняют отношения. Неизменность крупных контактов противоречит чисто равномерному масштабированию.}
